\documentclass[10pt,a4paper]{amsart}
\usepackage[margin=19mm]{geometry}
\usepackage{amsmath,amssymb,mathtools}
\usepackage[scr=rsfs,cal=euler]{mathalfa}
\usepackage{xfrac}
\usepackage[unicode,hidelinks,bookmarks=false]{hyperref}
\newcommand{\ie}{i.e.\ }
\newcommand{\cf}{cf.\ }
\newcommand{\resp}{respectively\ }
\newcommand{\Subsection}{\subsection}
\providecommand{\nobreakdash}{\nobreak\hbox{-}}
\setlength{\emergencystretch}{2em}
\multlinegap0pt
\usepackage{bm}
\usepackage{bbm}
\usepackage{dsfont}
\usepackage{xspace}
\usepackage{cancel}
\usepackage{mathtools}
\usepackage{amsthm}
\makeatletter
\@ifundefined{thm}{\newtheorem{thm}{Thm}}{}
\@ifundefined{lemma}{\newtheorem{lemma}[thm]{Lemma}}{}
\makeatother
\title[Square Functions and Variational Estimates for Ritt Operator]{Square Functions and Variational Estimates for Ritt Operators on $L^1$}
\author{Jennifer Hults, Karin Reinhold-Larsson}
\date{Source: arXiv:2507.07256v2 (CC BY 4.0)}
\begin{document}
\maketitle

\noindent\emph{Source and licence.} Square Functions and Variational Estimates for Ritt Operators on $L^1$. Version arXiv:2507.07256v2 (CC BY 4.0), distributed under the Creative Commons Attribution 4.0 International licence, as recorded in the arXiv licence metadata for this exact version and shown on the abstract page https://arxiv.org/abs/2507.07256v2. Full rights notice: licenses/arxiv-2507.07256-CC-BY-4.0.txt.\par\smallskip

\noindent\textit{Editorial scope.} Counted unit: the Lemma whose source label is \texttt{lem:basic} in the article (source0.tex lines 494--506, source label \texttt{lem:basic}), delivered together with the article's own complete original proof of it (source0.tex lines 509--562, one \texttt{proof} environment of 54 lines). No mathematics is written, repaired or re-derived: statement and proof are the article's own text line for line. Required context retained uncounted: none. Every definition, display and label of those lines is retained unchanged. Omitted from this package and not redistributed: 1--493, 507--508, 563--1495. Nothing inside a retained excerpt is omitted or altered: every retained body is delivered exactly as the article's lines are written, so no omission receipt of any kind accompanies this package. Citations: the retained excerpts carry no citation command at all, so no bibliography entry of the article is transcribed and no reference is redistributed. Reference adaptation: none was needed and none was made; every reference inside the delivered unit resolves inside it, so no unresolved reference or citation is produced. Printed slips kept literal: no printed word, symbol, hypothesis or number of the counted statement or of its proof was altered. Layout and class: the article's own class is \texttt{amsart} with packages including amsfonts, amsmath, amssymb, amsthm, comment, hyperref; the wrapper is the lane's pinned \texttt{amsart} editorial wrapper at 10pt on A4, which restates the theorem-like environments the retained text needs and adds the article's own macro definitions that the retained text needs (none), with extra wrapper packages (bm, bbm, dsfont, xspace, cancel, mathtools). The delivered numbering is the unit's own. Assets: no retained span contains a figure environment, a graphics-inclusion command, a PDF-page inclusion, a table or a TikZ picture; no image file is copied into this package, the package declares no asset dependency and no image file is redistributed with it. Licence: version arXiv:2507.07256v2 of this article is distributed under the Creative Commons Attribution 4.0 International licence, as recorded in the arXiv licence metadata for this exact version and shown on the abstract page https://arxiv.org/abs/2507.07256v2. A full rights notice is delivered at licenses/arxiv-2507.07256-CC-BY-4.0.txt. Characters: every retained excerpt is pure ASCII, so the delivered engine, XeLaTeX with its default Latin Modern text font, reports no missing character.
\begin{lemma}\label{lem:basic} 
Let $S=\{\Delta_n\}$ be a sequence of (finite) signed measures on $\mathbb Z$ whose Fourier transforms $\hat\Delta_n$ 
are twice continuously differentiable on $|t|\in (0,1)$. Let $T$ be a linear (doubly) power--bounded operator on a Banach function space, and
let
\begin{align*} 
A= & \int_{|t|<1/2} \frac 1{|t|} \Big( \sum_n |\hat \Delta_n(t)|^s \Big)^{1/s} dt,  \quad
& B=  \int_{|t|<1/2} |t| \Big( \sum_n |\hat \Delta''_n(t)|^s \Big)^{1/s} dt, \\
 C=& \sum_{k\neq 0} \frac 1{|k|} \Big( \sum_n |\hat \Delta_n(1/|2k|)|^s \Big)^{1/s},  \quad
 & D= \sum_{k\neq 0} \frac 1{|k|^2} \Big( \sum_n |\hat \Delta'_n(1/|2k|)|^s \Big)^{1/s}. 
\end{align*}
If $A, B, C$ and $D$ are all finite,  then, for any $f\in Y$,
$\Big\lVert  \Big( \sum_n |T_{\Delta_n}f|^s \Big)^{1/s} \Big\rVert \lesssim \lVert f\rVert.$
\end{lemma}

\begin{proof}
Without loss of generality, we assume $\sup_{n\ge 0} \lVert T^n \rVert = 1$ if $\cup_n \mbox{supp}(\Delta_n) \subset \mathbb N_0$ and  assume $\sup_{n\in \mathbb Z} \lVert T^n \rVert = 1$ otherwise.

\begin{align*}   \Big\lVert \Big( \sum_n |T_{\Delta_n}f|^s \Big)^{1/s} \Big\rVert
= & \Big\lVert \Big(\sum_n  \Big| \sum_k \Delta_n(k) T^kf\Big|^{s} \Big)^{1/s} \Big\rVert \\
\lesssim &   \Big\lVert   \Big( \sum_n  \Big|  \sum_{k\neq 0} \int_{|t|<1/|2k| } \hat \Delta_n(t) e(kt) dt  \,  T^kf \Big|^{s} \Big)^{1/s} \Big\rVert \\
 & +  \Big\lVert    \Big( \sum_n  \Big|  \sum_{k\neq 0}   \int_{1/|2k|<|t|<1/2} \hat \Delta_n(t) e(kt) dt  \,  T^kf \Big|^{s} \Big)^{1/s} \Big\rVert \\
& +     \Big( \sum_n  \Big|   \Delta_n(0)   \Big|^{s} \Big)^{1/s} \Big\lVert f \Big\rVert \\
 =  &  \mbox{I} + \mbox{II}+   \Big( \sum_n  \big|   \Delta_n(0)   \big|^{s} \Big)^{1/s} \Big\lVert f \Big\rVert.
\end{align*}



Applying the Stein--Minkowski integral inequality, we obtain the following estimate 
for the first term.
\begin{align*}
  \mbox{I} 
=  & \Big\lVert   \Big( \sum_n  \Big|  \sum_{k\neq 0} \int_{|t|<1/|2k|} \hat \Delta_n(t) e(kt) dt  \,  T^kf\Big|^{s} \Big)^{1/s} \Big\rVert \\
\le &  \lVert f\rVert \sum_{k\neq 0} \Big( \sum_n \Big|   \int_{|t|<1/ |2k| } \hat \Delta_n(t) e(kt) dt   \Big|^{s} \Big)^{1/s}\\
\le &   \lVert f\rVert \int_{|t|<1/2} \frac1{|t|}  \Big( \sum_n  \Big| \hat \Delta_n(t) \Big|^s \Big)^{1/s}  \, dt
= A \lVert f\rVert.
  \end{align*}
  Also 
\[ \Big( \sum_n  \big|   \Delta_n(0)   \big|^{s} \Big)^{1/s} \le  \int_{|t|<1/2} \Big( \sum_n  \Big| \hat \Delta_n(t) \Big|^s \Big)^{1/s}  \, dt
\le A.\]

For the second term, by the Stein--Minkowski integral inequality, we have
\begin{align*}
\mbox{II} \le & \Big\lVert   \Big( \sum_n  \Big|  \sum_{k\neq 0}   \int_{1/|2k| <|t|<1/2} \hat \Delta_n(t) e(kt) dt  \  T^kf \Big|^{s} \Big)^{1/s} \Big\rVert \\
\le & \Big\lVert  \sum_{k\neq 0}  \Big( \sum_n  \Big |  \int_{1/|2k|<|t|<1/2} \hat \Delta_n(t) e(kt) dt \Big |^s \Big)^{1/s} |T^kf | \Big\rVert\\
\le & \lVert f\rVert  \sum_{k\neq 0}  \Big( \sum_n  \Big |  \int_{1/|2k|<|t|<1/2} \hat \Delta_n(t) e(kt) dt \Big |^s \Big)^{1/s}.
\end{align*}

Since $\Delta_n(t), \Delta'_n(t)$ and $e(kt)$ are 1-periodic, $\Delta_n(-t)=\overline{\Delta_n(t)}$, and\\ $\Delta'_n(-t)=\overline{\Delta'_n(t)}$, we estimate
\begin{align}
\Big| \int_{1/|2k| <|t|<1/2} & \hat \Delta_n(t) e(kt) dt \Big| \le  
 \Big|\int_{1/|2k| <|t|<1/2} \hat \Delta'_n(t) \frac{e(kt)}{2\pi k} dt \Big| \nonumber \\
& + \Big|\frac{\hat \Delta_n(1/|2k| ) e(k/|2k| )}{2\pi k} - \frac{\hat \Delta_n(-1/|2k| ) e(-k/|2k| )}{2\pi k} \Big| \label{eq:1stderivative}\\   %eq2
\lesssim &  \Big| \int_{1/|2k| <|t|<1/2} \hat \Delta''_n(t) \frac{e(kt)}{4\pi^2 k^2} dt \Big| %\nonumber 
  + \Big| \frac{\hat \Delta_n(1/|2k| ) }{2\pi k} \Big| 
  +  \Big| \frac{\hat \Delta'_n(1/|2k| )}{4\pi^2 k^2}\Big|. 
 \nonumber
 \end{align}

By the Stein--Minkowski integral inequality,
\begin{align*}
\sum_{k\neq 0}   \Big( \sum_n  \Big |  \int_{1/|2k|<|t|<1/2} \hat \Delta''_n(t) \frac{e(kt)}{4\pi^2 k^2}  dt \Big |^s \Big)^{1/s} 
\lesssim  & \sum_{k\neq 0} \frac 1{k^2}  \int_{1/|2k|<|t|<1/2} \Big( \sum_n  \Big | \hat \Delta''_n(t) \Big |^s \Big)^{1/s} \ dt  \\
\lesssim & \int_{0<|t|<1/2}  |t| \,  \Big( \sum_n  \Big | \hat \Delta''_n(t) \Big |^s \Big)^{1/s}\, dt  = B .
\end{align*}

Thus, $\mbox{II}\lesssim (B+C+D) \lVert f\rVert$.

\end{proof}

\end{document}
